Данный раздел содержит описание предметной области, основные понятия, анализ текущего состояния процесса обращений граждан, постановку задачи, обзор аналогов и техническое задание на разработку программной системы.

\subsection{Описание предметной области}

Городская среда представляет собой сложную систему, в которой жители ежедневно взаимодействуют с объектами инфраструктуры и муниципальными службами. Город Киров~--- административный центр Кировской области с населением около 530~тысяч человек~--- включает четыре городских района (Ленинский, Октябрьский, Первомайский, Нововятский), каждый из которых обладает собственной инфраструктурой жилищно-коммунального хозяйства, объектами благоустройства и сетью ответственных городских служб.

Городская инфраструктура Кирова охватывает широкий спектр объектов. К основным категориям относятся:
\begin{itemize}
    \item дорожное покрытие и тротуары;
    \item уличное освещение;
    \item контейнерные площадки и объекты санитарной очистки;
    \item дворовые территории и объекты благоустройства.
\end{itemize}

Каждая категория объектов закреплена за определённой организацией или ведомством:
Каждая категория объектов закреплена за определённой организацией или ведомством. Обслуживание инфраструктуры распределено между муниципальными учреждениями, унитарными предприятиями, региональными операторами и управляющими компаниями.

Ключевой проблемой взаимодействия жителя с городской средой является отсутствие единого удобного канала для сообщения о дефектах инфраструктуры и контроля хода их устранения. Обслуживание городской инфраструктуры распределено между множеством организаций: за дорожное покрытие, уличное освещение, контейнерные площадки и иные объекты отвечают разные службы, каждая из которых имеет собственные каналы приёма обращений. Для жителя это означает необходимость самостоятельно определить, какая именно служба отвечает за устранение конкретной проблемы, найти её контактные данные и составить обращение по телефону или электронной почте.

Не менее существенные проблемы испытывают муниципальные службы на стороне диспетчера: диспетчеры не располагают единым инструментом для регистрации и ведения обращений, отсутствует визуализация проблем на карте города, определение ответственной службы по категории проблемы не автоматизировано, журнал действий оператора не ведётся.

На стороне жителя процесс подачи обращений также остаётся неавтоматизированным и сопряжён с рядом проблем:
\begin{itemize}
    \item житель вынужден самостоятельно определять, какая служба отвечает за устранение конкретной проблемы;
    \item контактная информация муниципальных служб рассредоточена по разделам сайта администрации города, ведомственным порталам и справочным ресурсам, актуальность найденных данных не гарантирована;
    \item отсутствует возможность оперативно зафиксировать проблему непосредственно в момент её обнаружения с мобильного устройства;
    \item после подачи обращения житель не получает подтверждения о его регистрации и не может отследить статус обработки;
    \item при отсутствии реакции приходится повторять весь процесс заново.
\end{itemize}

Ситуацию усугубляет отсутствие у жителей доступа к уже поданным обращениям: не зная о заявках других пользователей, разные люди сообщают об одной и той же проблеме. Механизм обнаружения дубликатов способен выявить повторные обращения, однако без мобильного приложения житель не может увидеть, что аналогичная проблема уже зафиксирована поблизости.

Таким образом, предметная область настоящей работы охватывает процесс фиксации проблем городской инфраструктуры жителями города Кирова и их обработку диспетчерами муниципальных служб.

\subsection{Понятия, относящиеся к предметной области}

Обращение~--- зафиксированная жителем неисправность или дефект городской инфраструктуры (яма на дороге, неисправный фонарь, переполненный мусорный контейнер и~т.\,п.), требующий реагирования со стороны ответственной городской службы. Обращение содержит текстовое описание, фотографию, категорию проблемы и географические координаты.

Категория проблемы~--- классификационный признак, позволяющий отнести обращение к одному из тематических направлений: дороги, освещение, благоустройство, ЖКХ, экология и другие. Категория определяет, какая городская служба является ответственной за устранение данного типа проблемы.

Ответственная служба~--- организация или подразделение, в компетенцию которого входит устранение проблем определённой категории. Для каждой службы в системе хранятся контактные данные: телефон, адрес, режим работы.

Справочник служб~--- структурированный перечень городских служб с их контактной информацией, заполняемый администратором системы на основании официальных открытых источников. Справочник обеспечивает привязку категорий проблем к ответственным исполнителям.

Статус обращения~--- текущее состояние обращения в системе. Жизненный цикл обращения включает следующие состояния: <<новое>> (обращение создано жителем), <<принято>> (диспетчер подтвердил обращение), <<в работе>> (обращение передано ответственной службе), <<решено>> (проблема устранена), <<отклонено>> (обращение признано нерелевантным).

История статусов~--- хронологическая запись всех изменений статуса обращения с указанием даты, времени и оператора, выполнившего изменение. История обеспечивает прозрачность обработки и возможность аудита.

Медиафайл обращения~--- фотография или иной файл, прикреплённый жителем к обращению при его создании через мобильное приложение. Медиафайлы хранятся во внешнем хранилище, а в базе данных сохраняются ссылки на них.

Геолокация~--- географические координаты (широта и долгота), привязанные к обращению и позволяющие отобразить его на интерактивной карте города. В мобильном приложении координаты определяются автоматически на основании данных GPS-модуля устройства с возможностью ручной корректировки пользователем.

Дубликат обращения~--- повторное обращение, описывающее ту же проблему, что и ранее поданное, и расположенное в непосредственной географической близости от него (в пределах заданного радиуса). Механизм обнаружения дубликатов реализован на стороне сервера.

Диспетчер~--- оператор системы, осуществляющий обработку поступивших обращений через веб-панель: просмотр на карте, определение ответственной службы, изменение статуса и контроль хода устранения проблемы.

Диспетчеризация~--- процесс обработки обращений, включающий их приём, классификацию, определение ответственной службы, изменение статуса и контроль результата. Диспетчеризация осуществляется посредством веб-панели.

Интерактивная карта~--- компонент программной системы, обеспечивающий визуальное отображение обращений в привязке к координатам на карте города. Карта используется как жителем в мобильном приложении, так и диспетчером в веб-панели.

Программный интерфейс (API)~--- набор серверных методов, обеспечивающих взаимодействие мобильного приложения и веб-панели с серверной частью системы.

Кроссплатформенная разработка~--- подход к созданию мобильных приложений, при котором единая кодовая база используется для сборки приложений под несколько платформ (Android и iOS). В рамках данной работы используется фреймворк React Native.

\subsection{Текущее состояние и проблемы взаимодействия жителей с городской средой}

В настоящее время в городе Кирове отсутствует единый цифровой канал для сообщения о проблемах городской инфраструктуры. Это создаёт ряд сложностей как для жителей, так и для муниципальных служб.

Со стороны жителя процесс подачи обращения выглядит следующим образом: обнаружив дефект инфраструктуры, гражданин вынужден самостоятельно определить, какая служба отвечает за его устранение, найти её контактные данные (рассредоточенные по разделам сайта администрации, ведомственным порталам и справочным ресурсам) и составить обращение по телефону или электронной почте. После отправки обращения житель не получает подтверждения о его регистрации и не может отследить статус обработки. При отсутствии реакции приходится повторять весь процесс заново.

К примеру, житель замечает крупную яму на дороге по пути на работу. Чтобы сообщить о ней, ему необходимо определить, кто отвечает за данный участок~--- управляющая компания, дорожная служба или администрация района. Далее нужно найти телефон или электронную почту этой организации, дозвониться или написать письмо. А если через неделю яма всё ещё на месте, приходится повторять весь процесс заново.

Не менее существенные проблемы испытывают муниципальные службы. Диспетчеры не располагают единым инструментом для регистрации и ведения обращений: заявки поступают по разным каналам (телефон, электронная почта, письменные заявления), фиксируются в разнородных журналах или не фиксируются вовсе. Отсутствует возможность визуализации проблем на карте города, что затрудняет оценку масштаба и приоритизацию. Определение ответственной службы по категории проблемы не автоматизировано~--- оператор вынужден самостоятельно искать контактные данные. Журнал действий оператора не ведётся, что исключает аудит и контроль качества обработки.

Ситуацию усугубляет отсутствие механизма обнаружения дубликатов: жители не видят обращения друг друга, поэтому одна и та же проблема может быть зафиксирована многократно разными заявителями. Это увеличивает нагрузку на службы и затрудняет оценку приоритетности обращений.

Текущий процесс сообщения жителя о проблеме городской инфраструктуры (до внедрения программной системы) представлен на контекстной диаграмме в нотации IDEF0 (рисунок~\ref{fig:current_asis_a0}).

\begin{figure}[H]
    \centering
    \includegraphics[width=0.8\linewidth]{current_a-0.png}
    \caption{Текущий процесс сообщения о проблеме городской инфраструктуры (уровень А-0)}
    \label{fig:current_asis_a0}
\end{figure}

Декомпозиция текущего процесса (рисунок~\ref{fig:current_asis_decomp}) включает следующие подпроцессы:
\begin{enumerate}
    \item поиск ответственной службы~--- житель самостоятельно определяет, какая организация отвечает за устранение проблемы;
    \item отправка обращения~--- житель связывается со службой по телефону или электронной почте;
    \item ожидание реакции и обработка~--- отсутствие обратной связи, невозможность отследить статус;
    \item повторная подача обращений~--- при отсутствии реакции процесс повторяется заново.
\end{enumerate}

\begin{figure}[H]
    \centering
    \includegraphics[width=0.8\linewidth]{current_a0.png}
    \caption{Декомпозиция текущего процесса сообщения о проблеме (уровень А0)}
    \label{fig:current_asis_decomp}
\end{figure}

Разработка программной системы, объединяющей мобильное приложение для жителей, серверную часть с базой данных и диспетчерскую веб-панель, позволит устранить описанные проблемы. Житель получит возможность зафиксировать проблему непосредственно на месте через мобильное приложение и отслеживать статус обращения. Диспетчер получит единый инструмент для обработки обращений на интерактивной карте с доступом к контактам ответственных служб, фильтрацией по категориям и статусам, механизмом ведения истории статусов и серверной проверкой дубликатов.

\subsection{Моделирование целевого процесса}

Процесс фиксации и диспетчеризации городских проблем в разрабатываемой программной системе представлен на контекстной диаграмме в нотации IDEF0 (рисунок~\ref{fig:context_a0}). На вход поступает проблема инфраструктуры, обнаруженная жителем. Управление осуществляется на основании законодательства РФ и справочника категорий проблем и служб. Механизмами являются: житель (через мобильное приложение), программная система фиксации событий и диспетчер (через веб-панель). На выходе~--- устранённая проблема и уведомление жителя о результате.

\begin{figure}[H]
    \centering
    \includegraphics[width=0.8\linewidth]{to-be_a-0.png}
    \caption{Контекстная диаграмма процесса фиксации и диспетчеризации городских проблем (уровень А-0)}
    \label{fig:context_a0}
\end{figure}

Декомпозиция процесса (рисунок~\ref{fig:decomp_a0}) включает четыре подпроцесса:
\begin{enumerate}
    \item фиксация проблемы жителем через мобильное приложение~--- житель выбирает категорию, прикрепляет фотографию, указывает местоположение; обращение передаётся на сервер;
    \item проверка на дубликаты~--- серверный алгоритм выполняет геопространственный поиск существующих обращений в радиусе 50~м; при обнаружении дубликата жителю предлагается подтвердить существующее обращение;
    \item диспетчеризация~--- диспетчер обрабатывает обращение через веб-панель: просматривает на карте, определяет ответственную службу, изменяет статус;
    \item контроль исполнения~--- отслеживание статуса обращения, фиксация результата, уведомление жителя через мобильное приложение.
\end{enumerate}

\begin{figure}[H]
    \centering
    \includegraphics[width=0.8\linewidth]{to_be_a0.png}
    \caption{Декомпозиция процесса фиксации и диспетчеризации (уровень А0)}
    \label{fig:decomp_a0}
\end{figure}

Подпроцесс <<Диспетчеризация>> детализирован на диаграмме в нотации IDEF3 (рисунок~\ref{fig:dispatcher_idef3}). Диаграмма описывает последовательность действий оператора при обработке обращения, поступившего от жителя:
\begin{enumerate}
    \item просмотр поступившего обращения~--- диспетчер открывает карточку в веб-панели, система отображает описание, фотографию (загруженную жителем), категорию и местоположение на карте;
    \item перекрёсток~J1~--- на основании просмотра принимается решение: <<обращение нерелевантно>> или <<обращение корректное>>;
    \item отклонение обращения~--- диспетчер указывает причину, статус изменяется на <<отклонено>>, житель получает уведомление;
    \item определение ответственной службы~--- система определяет службу по категории проблемы;
    \item ввод данных об обращении и редактирование~--- уточнение категории, описания, ввод комментария;
    \item передача обращения в работу~--- статус изменяется на <<в работе>>, житель получает уведомление;
    \item фиксация действий в журнале аудита~--- система сохраняет запись о действии с указанием даты, времени и оператора.
\end{enumerate}

\begin{figure}[H]
    \centering
    \includegraphics[width=0.7\linewidth]{to_be_dispatcher_idef3_decompose.png}
    \caption{Декомпозиция подпроцесса <<Диспетчеризация>> в нотации IDEF3}
    \label{fig:dispatcher_idef3}
\end{figure}

\subsection{Обзор аналогов}\label{subsec:analogs}

Для обоснования разработки собственной системы были рассмотрены существующие решения в области фиксации городских проблем и обработки обращений граждан. Анализ проводился с акцентом на полноту функционала: наличие мобильного приложения, кроссплатформенность, автоматическое определение геолокации, интеграцию с камерой устройства, уведомления, просмотр обращений на интерактивной карте и механизм дедупликации.

\subsubsection{Платформа обратной связи <<Госуслуги. Решаем вместе>>}

Платформа обратной связи (ПОС)~--- федеральный сервис, интегрированный с порталом Госуслуг, предназначенный для подачи обращений граждан в органы власти. Помимо веб-сервиса, в магазинах приложений доступно отдельное приложение <<Госуслуги Решаем вместе>> (рисунок~\ref{fig:analog-gosuslugi}). Достоинства: федеральный охват, включая город Киров; наличие мобильного приложения; поддержка фотографий и карты обращений. Недостатки: нет механизма дедупликации; автоматическое определение GPS-координат не является приоритетным сценарием; маршрутизация к конкретной ответственной службе не реализована.

\begin{figure}[H]
    \centering
    \includegraphics[width=0.5\linewidth]{gosuslugi-reshaem-vmeste.jpg}
    \caption{Мобильное приложение <<Госуслуги. Решаем вместе>>}
    \label{fig:analog-gosuslugi}
\end{figure}

\subsubsection{Портал <<Наш город>> (Москва)}

<<Наш город>>~--- московский портал для подачи жалоб на состояние городской инфраструктуры с развитой системой категоризации, механизмом верификации и мобильным приложением (рисунок~\ref{fig:analog-nashgorod}). Поддерживается фотофиксация, автоопределение геолокации, push-уведомления. Недостатки: доступен исключительно для жителей Москвы; автоматическая дедупликация по радиусу отсутствует.

\begin{figure}[H]
    \centering
    \includegraphics[width=0.5\linewidth]{nash-gorod-moskva-user.jpg}
    \caption{Мобильное приложение <<Наш город>> (Москва)}
    \label{fig:analog-nashgorod}
\end{figure}

\subsubsection{Добродел (Московская область)}

<<Добродел>>~--- портал обращений жителей Московской области с механизмом <<коллективного сообщения>> и кнопкой <<Поддержать>> (рисунок~\ref{fig:analog-dobrodel}). Мобильное приложение поддерживает фотофиксацию, автоопределение геолокации и push-уведомления. Недостатки: ограничен территорией Московской области; автоматическая дедупликация по радиусу не реализована; интерфейс перегружен дополнительными функциями.

\begin{figure}[H]
    \centering
    \includegraphics[width=0.5\linewidth]{dobrodel-user.jpg}
    \caption{Мобильное приложение <<Добродел>>}
    \label{fig:analog-dobrodel}
\end{figure}

\subsubsection{<<Наш Санкт-Петербург>>}

<<Наш Санкт-Петербург>>~--- региональный сервис для фиксации городских проблем в Санкт-Петербурге (рисунок~\ref{fig:analog-nashspb}). Мобильное приложение поддерживает фотофиксацию, указание местоположения, push-уведомления. Недостатки: ограничен территорией Санкт-Петербурга; автоматическая дедупликация не реализована.

\begin{figure}[H]
    \centering
    \includegraphics[width=0.5\linewidth]{nash-spb-user.jpg}
    \caption{Мобильное приложение <<Наш Санкт-Петербург>>}
    \label{fig:analog-nashspb}
\end{figure}

\subsubsection{FixMyStreet}

FixMyStreet~--- британская платформа с открытым исходным кодом для фиксации городских проблем (рисунок~\ref{fig:analog-fixmystreet}). Достоинства: автоматическое определение геолокации, интерактивная карта, автоматическая маршрутизация, открытый исходный код. Недостатки: ориентирован на Великобританию; push-уведомления не реализованы; адаптация под российские реалии не выполнена.

\begin{figure}[H]
    \centering
    \includegraphics[width=0.5\linewidth]{fixmystreet-user.png}
    \caption{Мобильное приложение FixMyStreet}
    \label{fig:analog-fixmystreet}
\end{figure}

\subsubsection{Выводы по результатам анализа}

Проведённый сравнительный анализ выявил ряд ограничений, обосновывающих разработку собственной системы:
\begin{itemize}
    \item среди рассмотренных аналогов отсутствует решение, применимое к городу Кирову и обеспечивающее полный набор функций: автоматическое определение геолокации, отслеживание статуса, просмотр обращений на карте и механизм дедупликации;
    \item платформа ПОС, применимая к Кирову, не реализует дедупликацию и автоматическое определение GPS-координат;
    \item региональные решения (<<Наш город>>, <<Добродел>>, <<Наш Санкт-Петербург>>) территориально ограничены;
    \item FixMyStreet не адаптирован к российским реалиям и не поддерживает push-уведомления;
    \item ни один из рассмотренных аналогов не реализует автоматическую дедупликацию обращений по заданному радиусу.
\end{itemize}

\subsection{Постановка задачи}

Необходимо разработать программную систему фиксации городских событий для жителей города Кирова, включающую три компонента:
\begin{enumerate}
    \item серверную часть~--- программный интерфейс (REST API) и реляционную базу данных PostgreSQL с расширением PostGIS;
    \item диспетчерскую веб-панель~--- автоматизированное рабочее место оператора с интерактивной картой обращений;
    \item кроссплатформенное мобильное приложение на базе React Native для подачи обращений жителями.
\end{enumerate}

Основные функциональные модули системы:
\begin{enumerate}
    \item регистрация и аутентификация пользователей;
    \item создание обращения о проблеме городской инфраструктуры~--- фото, категория, описание, автоматическое определение геолокации;
    \item просмотр обращений на интерактивной карте с фильтрацией и режимом тепловой карты;
    \item обработка обращений диспетчером~--- просмотр на карте, определение ответственной службы, изменение статуса;
    \item автоматическое обнаружение дубликатов обращений в радиусе 50~м;
    \item маршрутизация обращений к ответственным службам по категории проблемы;
    \item просмотр списка собственных обращений и отслеживание их статуса;
    \item уведомление пользователя об изменениях статуса обращения;
    \item справочники категорий проблем и муниципальных служб.
\end{enumerate}

Диаграмма прецедентов мобильного приложения представлена на рисунке~\ref{fig:use_case_mobile}.

\begin{figure}[H]
    \centering
    \includegraphics[width=0.8\textwidth]{use-case-mobile.png}
    \caption{Диаграмма прецедентов мобильного приложения}
    \label{fig:use_case_mobile}
\end{figure}
